\section{Discussion}
\label{sec:discussion}

\subsection{A Unified Flow without an Integrated Claim}

The three components address one recurring security-learning problem: useful
evidence is incomplete, arrives at different times, and has different costs.
Static and dynamic scores allocate evidence to an online decision. Active
learning allocates a label budget to model maintenance. TreeSHAP diagnostics
allocate confidence to an explanation ranking used in that maintenance process.
This flow is scientifically useful because it identifies the interfaces a future
system must validate: paired channel identity, acquisition lineage, and
explanation checks.

The flow is not itself an empirical interaction. MalBehavD-V1 does not contain
the EMBER active-learning artifact, and EMBER does not contain the ordered
dynamic observation required for fusion. Consequently, the study cannot answer
whether explanation-guided retraining improves the fused detector, whether PID
changes which samples should be queried, or whether one common threshold reduces
end-to-end alert cost. A factorial experiment on one paired cohort remains
necessary. Preserving this boundary makes the component results easier to audit
and prevents a conceptual architecture from becoming a fabricated system score.

\subsection{What the PID Translation Establishes}

The control analogy provides a disciplined set of nested hypotheses. P asks
whether current detector disagreement should alter the mixture. I asks whether
persistent signed disagreement supplies additional information. D asks whether a
change in disagreement across event prefixes predicts useful adjustment. The
empirical result supports the first hypothesis and rejects the need for the
latter two under the measured horizon.

This outcome is plausible. The static score is constant within a sequence, while
the dynamic score updates at short API prefixes. If successive prefixes change
only a few events, the derivative term can be sparse or noisy. If one channel is
systematically higher for a source, the integral may accumulate collection bias
rather than sample-specific reliability. A proportional mapping can capture the
dominant behavior with fewer states and lower implementation risk. Adaptive P's
FPR and Brier improvements over Fixed make it a stronger current candidate than
full PID for horizon 1--20.

The result also clarifies what PID cannot do. It cannot reveal behavior that the
sandbox never observes. It cannot turn event count into elapsed time. It cannot
guarantee calibration or stability from control theory alone. Sleep-evasion
resistance would require real timestamped traces, controlled delayed behavior,
and a specified observation policy. End-to-end latency would require separate
measurement of queueing, sandbox startup, observation, feature extraction,
channel inference, and fusion. The small timestamp pilot begins data-source
validation but does not satisfy those requirements.

Learned stacking remains an important baseline. Full PID has comparable AUC and
lower FPR at its own operating point, but also a small TPR reduction. Average has
the best Brier and slightly higher AUC than the adaptive family. These results
show that interpretability and operating-point objectives, rather than a tiny
AUC difference, must motivate an adaptive controller. External validation at a
common target TPR should decide whether Adaptive P's lower FPR transfers.

\subsection{What the Acquisition Result Means}

Directional entropy adds a small and repeatable learning-curve improvement over
Random and a smaller improvement over Uncertainty in four of five folds. It
distinguishes two probability-equivalent cases: weak evidence and stronger
opposing signed masses. Such conflict can identify a model region where a new
label is useful. The algebra in Equation~(\ref{eq:h-identity}) narrows this
interpretation: $H$ is determined by the total raw-score contribution and the
$L_1$ attribution norm. A direct margin-plus-norm baseline may recover part of
the effect and should be included in an extension.

No acquisition policy dominates every objective. $U+H$ leads AULC, $U+V$
reaches the target earliest, \xfull balances signals and saves 92 labels
relative to Random, and Uncertainty remains strongest in two restricted-FPR
summaries. Policy selection should therefore follow the maintenance objective.
An organization prioritizing average model quality during iterative labeling may
prefer $U+H$; one using the retained validation target may prefer $U+V$; one
requiring low computation and strong low-FPR behavior may retain Uncertainty.

The label result has an explicit accounting boundary. It counts labels added to
training. Validation labels, initial dataset creation, the full-pool reference,
analyst abstention, and disagreement resolution remain outside the number. A
15.8\% label-count reduction is not a measured 15.8\% reduction in labor.

Batch redundancy is an unresolved alternative explanation. Selecting the top
100 pointwise scores can query many objects from one region. Uncertainty plus
k-center selection would test whether simple diversity reproduces the gain. A
calibrated committee would test whether ensemble disagreement explains the
benefit. Robust quantile or rank normalization would test sensitivity to extreme
$H$ or $V$ values that compress min--max scaling. These settings must be frozen
within a nested design, not chosen by outer-test AULC.

\subsection{What the TreeSHAP Audit Establishes}

Top-ranked replacement damages model ranking much more than uniform random
replacement under both training means and medians. This is strong evidence of
relative model reliance under the defined feature-space intervention. It is
particularly useful as an internal release check: a TreeSHAP ranking that fails
to exceed equal-count random deletion should not be presented as a meaningful
ordering of influential coordinates.

The deletion gap is not a causal certificate. Tree ensembles use a subset of
2,381 coordinates, so random masks often strike inactive features. Top-ranked
coordinates may also be farther from the replacement baseline, creating a larger
standardized displacement and greater off-manifold shift. Stronger controls
should sample only used features, match displacement magnitude, replace
semantically linked PE groups, and use conditional backgrounds. Reconstructing
valid PE binaries would answer a different and substantially harder question.

Moderate stability gives a complementary warning. Explanations change less
across model seeds in one fold than across folds, showing that data composition
matters. A global top-feature list should therefore be versioned with the model,
feature schema, and training cohort. For active learning, round-specific change
is expected because the model is retrained, but abrupt attribution drift should
be logged and investigated.

Additivity, EC, and CF remain diagnostics of numerical or geometric properties.
A common permutation of attribution columns can preserve sums, cosines, norms,
and centroid separation while giving effects the wrong feature identities.
Combining such heterogeneous quantities into an unvalidated Overall Fidelity
Score would hide this weakness. Separate reporting is more informative and more
defensible.

\subsection{Operational Interpretation}

The most immediate use is a controlled model-development workflow. A candidate
static model scores a pool; a chosen acquisition rule prioritizes objects for an
authorized analyst; query/model lineage is retained; and deletion/stability
checks are run before release. Separately, when a paired dynamic sequence is
available, Adaptive P can be compared against Average, Fixed, and Stacking at a
locked operating objective. The evidence does not yet authorize automated
remediation or deployment claims.

Scaling introduces additional costs. TreeSHAP over a large candidate pool can be
reduced by first filtering a probability band, but that changes the acquisition
population and requires reevaluation. Dynamic collection dominates arithmetic
fusion cost, so optimizing the PID calculation alone will not solve system
latency. These practical distinctions follow directly from the unified evidence-
allocation view: computation, labels, observation time, and alert errors are
different resources and should be measured separately.

\subsection{Evidence Needed for a Stronger Journal Claim}

The highest-priority experiment is an external or temporal paired cohort whose
method and thresholds are frozen before access. Adaptive P and robust PID v2
should be preregistered with a common-TPR comparison and noninferiority or
superiority margins. Timestamped event sequences should support separate early-
detection and wall-clock analyses. The same cohort should then receive a
factorial acquisition--fusion evaluation if interaction is claimed.

For RQ2/RQ3, the strongest additions are uncertainty plus diversity, a calibrated
committee, weight/normalization sensitivity, family- or time-aware splits,
used-feature and magnitude-matched deletion controls, conditional/group masking,
and a second explainer. These experiments address plausible alternatives to the
current effects and determine how far they generalize beyond the audited subset.
